Building a physical model for a process outside the libraries of modeling tools requires an
engineer to collect equations scattered across many documents and to assemble them into a model.
This study addresses one step of automating this work: given a description of the desired model and
its input and output variables, present the equations likely to serve in the model, drawn from a
database of equations extracted from the literature. Because an equation that contains a variable
other than the inputs and outputs is useful only together with an equation that determines that
variable, we propose a two-stage retrieval method that judges each candidate by its relation to
other equations. The first stage narrows the database by text similarity and variable overlap; the
second reranks the candidates with a multilayer perceptron that reads seven query--equation
features and three set-aware features. On a database of \NEq{} equations and \NCases{} cases split
by source, the method raised Recall@20, the fraction of the ground-truth equations among the 20
presented, from \RecABase{} for classical retrieval to \RecAProp, in all 10 splits. Ablation and
permutation importance show that the features reading the variables shared among equations cause
most of the gain, whereas topical similarity adds little. Two large language models that read the
whole database without training reached \RecAOpus{} and \RecAHaiku; the proposed method beat both
in every split. The larger model came close on models drawn from a single source and found
equations that share few variables with the query, at \CostPerCallOpus{} US dollars per query and
with answers that varied between runs.
